\subsection{Videovigilancia y Seguridad en el Entorno Retail}

En el entorno retail, los sistemas tradicionales de Circuito Cerrado de
Televisión (CCTV) han operado desde su nacimiento como un mecanismo pasivo,
cuya eficacia depende casi completamente de operadores humanos encargados de
supervisar las imágenes en tiempo real. Sin embargo, este último punto es la
principal desventaja de estos métodos pues, el operador humano se ve
completamente saturado por la cantidad de pantallas, ángulos y situaciones
distintas a las que debe prestar atención, sumado a esto están factores
como la falta de experiencia, aburrimiento, distracciones, sesgos del
operador, entre otros.

Frente a estas desventajas, la videovigilancia inteligente (\textit{Intelligent Video Surveillance -- IVS) se posiciona como un
paradigma de vigilancia activo basado en modelos de Deep Learning que transforman
los flujos de imágenes de diversas fuentes en información estructurada,
con anotaciones visuales y alertas en tiempo real. En el contexto retail y de este
trabajo, nos enfocamos en aquellas soluciones planteadas para la detección precisa
y temprana de hurtos en tiendas peruanas minoristas --modalidad conocida en la literatura
como \textit{shoplifting}, con miras a ofrecer una solución capaz de mitigar la merma
desconocida, es decir aquella que proviene directamente de hurtos.

La ventaja principal de la IVS está en la capacidad de operar continuamente,
observando diversas fuentes de datos simultáneas y filtrando solamente los segmentos de
imagen relacionados a actividades de hurto. Gracias al crecimiento del deep learning
y la literatura en este campo, podemos encontrar soluciones con diversos enfoques, enfocadas en
la detección directa en tiempo real del hurto con modelos espaciales como \cite{Aguero, Salah},
también aquellos que plantean que el hurto no es una actividad aislada en el tiempo, sino que
se da en 4 fases \cite{martinez, cinthula} por lo cual es necesario utilizar modelos
espacio-temporales, asimismo autores que combinan modelos espaciales con sensores de proximidad, IR
y de peso para determinar si se ha dado realmente un hurto \cite{Pranav}, por último tenemos las
soluciones orientadas a proteger la privacidad sin sacrificar rendimiento en la detección, propuesta por
\cite{Poselift} quien también aporta un dataset \textit{Poselift} específicamente realizado con este propósito.

\subsection{Modelos Espaciales y Espacio-Temporales}

Los modelos espaciales son aquellos que parten de una CNN capaz de aprender \textit{features} espaciales y
reconocer patrones de crimen a través de estos \cite{Chioma}, este tipo de modelos resultan prevalentes en la
literatura, especialmente en los trabajos que utilizan la arquitectura YOLO, un CNN \texit{single-stage}, que
procesa las imágenes en una sola pasada, predice la generación de \textit{bounding boxes} --cajas delimitadoras del
objeto detectado--, y categoriza el objeto detectado en tiempo real. El trabajo realizado por \cite{Aguero}, también
en el contexto retail peruano, utiliza YOLOv8 para detectar en tiempo real el hurto y señalar este dentro de una
bounding box en tiempo real. También vemos como Pranav et al (\cite{m_deep_2023}) suma sensores IR, de proximidad y de
peso a este modelo para realizar una detección más precisa del hurto apoyado de varios fuentes de confirmación. Por
último es importante mencionar que este tipo de modelos son mucho menos costosos computacionalmente, capaces de correr
en \textit{edge devices} como NVIDIA Jetson (\cite{hameed}), Raspberri Pi 3 (\cite{m_deep_2023}) y hardware limitado
o antiguo (\cite{Aguero}).

Sin embargo, diversos autores plantean que el crimen debe analizarse como una secuencia de etapas conductuales a lo largo
del tiempo, es aquí donde entran los modelos espacio-temporales con memoria recurrente. Las RNN (\textit{recurrent neural network)
son redes neuronales que introducen un circuito interno de retroalimentación que actúa como una memoria del contexto pasado,
sin embargo sufren del problema de desvanecimiento al procesar secuencias de prolongadas, como si estuvieran olvidando. Para
solucionar esto, se introducen las redes LSTM (\textit{Long Short-Term Memor}) las cuales utilizan una celda de memoria controlada
por tres mecanismos de compuerta, la \textit{forget gate} que utiliza una función sigmoide para decidir que elementos de la celda
deben olvidarse, la \textit{input gate} decide qué información nueva debe guardarse en el estado actual de la celda, y la
\textit{output gate} decide que parte del estado de la celda actualizada se enviará hacia el siguiente paso de la secuencia. Este
circuito de memoria es especialmente útil en la detección de hurtos, pues como propone Martinez-Mascorro \cite{martinez}, el crimen
se da durante cuatro etapas para las cuales plantea el framework PCB (\textit{Pre-Crime Behavior}).